%% tags: [best and worst case]
%% source: 2023-fa-midterm_01
\begin{prob}[(2 points)]
    The code below takes in two lists, each containing $n$ integers, and
    determines if the any number in the second list appears at least $n/2$
    times in the first list.

    \inputminted[lastline=12]{python}{code.py}

    \begin{subprobset}

        \begin{subprob}
            What is the \textbf{best case} time complexity of this code as a function of $n$?
            State your answer using asymptotic notation.

            \inlineresponsebox{$\Theta(n)$}

        \end{subprob}

        \begin{subprob}

            What is the \textbf{worst case} time complexity of this code as a function of $n$?
            State your answer using asymptotic notation.

            \inlineresponsebox{$\Theta(n^2)$}

        \end{subprob}

    \end{subprobset}

    \begin{soln}
        The best case is when the first element of the second list appears at least $n/2$ times
        in the first list. In this case, the code will return \mintinline{python}{True} after iterating $n/2$ times
        through the inner loop, taking $\Theta(n)$ time total.

        In the worst case, there is no element in the second list that appears
        $n/2$ times in the first. We make all $n$ iterations of the outer loop,
        and during each of the outer iterations, make $n$ iterations of the
        inner loop, for a total of $\Theta(n^2)$ time.
    \end{soln}

\end{prob}
